% TOP_PROBLEM: {"id": 1862, "title": "Twin Prime Conjecture", "category_id": 1, "category": {"id": 1, "name": "number_theory", "display_name": "Number Theory", "description": "Properties of integers, prime numbers, Diophantine equations.", "slug": "number-theory", "order_index": 1, "created_at": "2026-07-31T15:26:25.670Z"}, "status": "open"}
% FIRST_POSTED: null
% ENTRY_KIND: statement_only
% Imported from: batch08.tex; UnsolvedMath ID 1862
\documentclass[11pt]{article}
\usepackage[margin=25mm]{geometry}
\usepackage{fontspec}
\setmainfont{Latin Modern Roman}
\usepackage{amsmath,amssymb,mathtools}
\usepackage{unicode-math}
\setmathfont{Latin Modern Math}
\usepackage{xurl,hyperref}
\hypersetup{colorlinks=true,urlcolor=blue}
\setcounter{secnumdepth}{0}
\setlength{\parindent}{0pt}
\setlength{\parskip}{5pt}
\setlength{\emergencystretch}{3em}
\newcommand{\sref}[2]{\hyperlink{source:#1:#2}{[S#2]}}
\newcommand{\cataloguescope}[1]{\paragraph{Recorded target.}#1}
\begin{document}
% UnsolvedMath ID 1862; source catalogue code GUY-A8b
\setcounter{section}{1861}
\section{Twin Prime Conjecture}\label{1862}
{\small\sffamily UnsolvedMath ID 1862 · Number Theory}
\subsection{Definitions and mathematical statement}

\paragraph{Shared notation.}$\mathbb N=\{1,2,\ldots\}$, and $\mathbb Z,\mathbb Q,\mathbb R,\mathbb C$ denote the integers, rationals, reals and complex numbers. Any other conventions are specified locally.


Two primes form a twin-prime pair if they have the form
\(p,p+2\). Count each pair by its smaller prime:
\[
                   T(x)=\#\{p\le x:p\text{ and }p+2\text{ are prime}\}.
\]
\paragraph{Statement recorded in the supplied source.}
Are there infinitely many twin-prime pairs? Equivalently,
does every bound \(B\in\mathbb N\) admit a prime \(p>B\) such that
\(p+2\) is also prime, or does \(T(x)\to\infty\) as \(x\to\infty\)? \sref{1862}{2} \sref{1862}{3}
\subsection{Short English statement}
Are there infinitely many pairs of primes separated by exactly two?
\subsection{Sources}
\begin{description}
\item[\hypertarget{source:1862:1}{S1}] ULAM.AI and the UnsolvedMath Contributors, \emph{UnsolvedMath: A Curated Collection of Open Mathematics Problems}, record 1862 (GUY-A8b). CC BY 4.0. \url{https://huggingface.co/datasets/ulamai/UnsolvedMath}.
\item[\hypertarget{source:1862:2}{S2}] James Maynard, \emph{Small gaps between primes}, Annals of Mathematics 181 (2015), 383--413; arXiv:1311.4600, Section 1, admissibility and the qualitative Prime $k$-tuples Conjecture. \url{https://arxiv.org/abs/1311.4600}.
\item[\hypertarget{source:1862:3}{S3}] D. H. J. Polymath, \emph{Variants of the Selberg sieve, and bounded intervals containing many primes}, Research in the Mathematical Sciences 1 (2014), Article 12; arXiv:1407.4897, abstract and Section 1, the fixed-gap twin-prime problem and partial bounded-gap results. \url{https://arxiv.org/abs/1407.4897}.
\end{description}
\label{end:1862}
\end{document}
